\chapter{Navigation Integration and Validation}
\label{ch:nav-integration}

% TODO: this chapter's final section (6.5) reports a QUALITATIVE
% end-to-end result ("working almost perfectly") without a fresh
% quantitative monitor_goal_timing.py re-run post-readMode-fix. If you
% can still get lab time on the QCar, re-running that tool and
% dropping the numbers into Table 6.2 (currently a TODO) would close
% the one open loop in this thesis's results and make Ch. 7 stronger.

This chapter describes bringing the full Nav2 navigation stack up
against the physical QCar, building on the bridge of
Chapter~\ref{ch:hardware-bridge} and the calibration of
Chapter~\ref{ch:calibration}. Section~\ref{sec:nav2-porting} covers
porting the simulation-tuned Nav2 configuration to hardware.
Section~\ref{sec:mppi-critic} describes a custom MPPI critic plugin
developed for this vehicle. Section~\ref{sec:phased-integration}
reports the phased bring-up used to de-risk the integration.
Section~\ref{sec:display-lag} presents a detailed timing
investigation into an apparent localisation-display artefact, and
Section~\ref{sec:final-validation} reports the resulting end-to-end
validation.

\section{Porting the Nav2 Configuration to Hardware}
\label{sec:nav2-porting}

The Nav2 parameter set developed and tuned in simulation
(Chapter~\ref{ch:sim-stack}) required a small, mechanical set of
changes to run against physical hardware rather than Gazebo: every
occurrence of \texttt{use\_sim\_time} (13 in total, across the
costmap, planner, controller, and AMCL parameter blocks) was switched
from \texttt{true} to \texttt{false}, since the physical robot has no
simulation clock to synchronise against. Both mapping and navigation
launch files were given a real-hardware branch, gated on
\texttt{launch\_sim}, alongside the existing Gazebo branch, with
real hardware as the default and simulation opt-in
(\texttt{launch\_sim:=true}) -- a deliberate inversion of which mode
is the default, made once hardware became the project's active
development target, while keeping the simulation path fully intact
for future regression testing rather than removing it.

The AMCL \texttt{recovery\_alpha\_fast}/\texttt{recovery\_alpha\_slow}
fix from Section~\ref{sec:ruled-out} was carried forward
proactively into the hardware configuration at this point, rather
than left to be independently rediscovered against hardware, since
the underlying cause -- Section~\ref{sec:amcl-math}'s augmented-MCL
random-particle-injection mechanism being permanently disabled at
$\alpha_{\text{slow}} = \alpha_{\text{fast}} = 0$, leaving the filter
with no way to recover from an overconfident convergence once real
sensor noise (which a clean simulation rarely exercises) triggers one
-- is a property of AMCL's own resampling algorithm, not specific to
any one deployment target.

\section{Custom MPPI Critic: \texttt{EarlyCommitCritic}}
\label{sec:mppi-critic}

Early navigation testing showed a specific weakness in the stock MPPI
critic set (Section~\ref{sec:nav2}) for this vehicle: the standard
path-following critics evaluate a sampled trajectory's alignment and
progress against the path as a whole, which can leave the controller
driving further into an approaching curve before it begins reacting
to it, rather than beginning the turn as early as physically sensible
-- a behaviour more forgiving for a differential-drive robot's tight
turning radius than for an Ackermann-steered vehicle with the
finite minimum turning radius of Table~\ref{tab:qcar-geometry}.

A custom MPPI critic plugin, \texttt{EarlyCommitCritic}, was
developed and built into the navigation package to address this: it
scores only the first \texttt{early\_time\_steps} of each sampled
trajectory rollout against the bearing to a near-term path point,
with no distance or angle gating, so that MPPI is pushed to begin
turning toward the path immediately rather than waiting until deeper
into the curve. It is gated by an \texttt{active\_path\_points}
parameter so that it stops influencing the controller once the
vehicle has made genuine progress along the route, avoiding
interference with the stock critics' normal path-tracking behaviour
during straight segments. A later direction-awareness refinement
extended the critic to handle reversing and K-turn path segments
correctly, where a naive ``always turn toward the next point''
bearing calculation would otherwise fight the vehicle's intended
direction of travel.

Building this plugin required matching \texttt{xtensor}/\texttt{xsimd}
compile definitions exactly against Nav2's own MPPI controller build
-- a mismatch here does not produce a compile-time error, but a
silent ABI incompatibility that crashes the navigation container on
the first control cycle that touches vectorised critic math, making
it a notably unforgiving class of build-configuration bug to diagnose
from its symptom alone.

Listing~\ref{lst:early-commit-critic} shows the critic's
\texttt{score()} method in full. It computes the bearing from each
sampled trajectory's own starting position to a near-term path target
(\texttt{offset\_from\_furthest} points ahead of the furthest point
the robot has already reached), then, for only the first
\texttt{early\_time\_steps} of that trajectory, accumulates the
angular distance between the trajectory's actual yaw and that
bearing -- scoring every sampled rollout independently, not just the
robot's single current pose. The early return on
\texttt{active\_path\_points} is what confines this to the start of a
path, and the \texttt{forward\_preference\_}-gated minimum of the
forward- and reverse-bearing angular distances is the
direction-awareness refinement mentioned above, letting a trajectory
be scored against whichever of the bearing or its reverse it is
actually closer to, rather than penalising a legitimate reversing
manoeuvre.

\begin{lstlisting}[language=C++,caption={\texttt{EarlyCommitCritic::score()} (\texttt{src/critics/early\_commit\_critic.cpp}).},label={lst:early-commit-critic}]
void EarlyCommitCritic::score(CriticData & data)
{
  if (!enabled_ || data.path.x.shape(0) < 2) { return; }
  utils::setPathFurthestPointIfNotSet(data);

  // Only relevant at the start of a path - once the robot has made
  // real progress, let the normal path/goal critics take over.
  if (*data.furthest_reached_path_point >= active_path_points_) { return; }

  const size_t path_size = data.path.x.shape(0) - 1;
  const size_t offseted_idx = std::min(
    *data.furthest_reached_path_point + offset_from_furthest_, path_size);
  const float target_x = data.path.x(offseted_idx);
  const float target_y = data.path.y(offseted_idx);

  const size_t time_steps = data.trajectories.x.shape(1);
  const size_t early_steps = std::min(early_time_steps_, time_steps);
  if (early_steps == 0) { return; }

  const xt::xtensor<float, 1> x0 = xt::view(data.trajectories.x, xt::all(), 0);
  const xt::xtensor<float, 1> y0 = xt::view(data.trajectories.y, xt::all(), 0);
  const xt::xtensor<float, 1> bearing = xt::atan2(target_y - y0, target_x - x0);
  const xt::xtensor<float, 1> reverse_bearing = bearing + static_cast<float>(M_PI);

  xt::xtensor<float, 1> sum_diff = xt::zeros<float>({data.trajectories.x.shape(0)});
  for (size_t t = 0; t < early_steps; ++t) {
    const xt::xtensor<float, 1> yaw_t = xt::view(data.trajectories.yaws, xt::all(), t);
    const xt::xtensor<float, 1> fwd_diff =
      xt::abs(utils::shortest_angular_distance(yaw_t, bearing));
    if (forward_preference_) {
      sum_diff += fwd_diff;
    } else {
      const xt::xtensor<float, 1> rev_diff =
        xt::abs(utils::shortest_angular_distance(yaw_t, reverse_bearing));
      sum_diff += xt::minimum(fwd_diff, rev_diff);
    }
  }
  const xt::xtensor<float, 1> cost = sum_diff / static_cast<float>(early_steps);
  data.costs += xt::pow(cost * weight_, power_);
}
\end{lstlisting}

\paragraph{How this cost actually changes which trajectory wins.}
Connecting this directly to Section~\ref{sec:mppi-math}'s softmax
weighting (Equation~\ref{eq:mppi-weight}): \texttt{data.costs} is the
same per-trajectory accumulator every enabled critic adds into, so
this critic's contribution
$\left(\overline{\Delta\theta} \cdot \texttt{weight\_}\right)^{\texttt{power\_}}$
(the early-steps-averaged angular difference $\overline{\Delta\theta}$
from the listing's \texttt{cost} variable, scaled and exponentiated by
its own configured weight and power) is simply summed alongside the
stock critics' own contributions into $S(\tau^{(k)})$ before
Equation~\ref{eq:mppi-weight}'s exponential weighting is applied.
Concretely: of two sampled trajectories approaching the same curve,
one that begins turning toward the near-term path point immediately
accrues a small $\overline{\Delta\theta}$ over its first
\texttt{early\_time\_steps}, while one that continues driving straight
before reacting accrues a much larger one -- inflating that second
trajectory's total cost $S(\tau^{(k)})$, which Equation~\ref{eq:mppi-weight}
translates into an \emph{exponentially} smaller importance weight
$w^{(k)}$ relative to the early-turning trajectory. Because the
executed control is the weight-averaged sum over all $K$ sampled
trajectories, systematically suppressing the weight of every
late-turning sample -- rather than merely marking one specific
trajectory as disqualified -- is what actually shifts the controller's
averaged output toward beginning the turn earlier, without ever
needing to identify or select a single ``best'' trajectory explicitly.

\section{Phased Hardware Integration}
\label{sec:phased-integration}

Given the number of independent subsystems involved (TCP bridge,
odometry, SLAM, localisation, planning, control), integration
proceeded in three deliberately staged phases rather than attempting
the full stack at once, so that a failure at any stage could be
attributed to a bounded set of newly-added components rather than the
entire system:

\begin{description}
  \item[Phase 1 -- Teleoperation.] Bridge (Chapter~\ref{ch:hardware-bridge})
    plus manual \texttt{/cmd\_vel} teleoperation only, with no
    odometry, mapping, or navigation involved. Confirmed the motor and
    steering command path worked correctly end-to-end before any more
    complex subsystem was layered on top.
  \item[Phase 2 -- SLAM.] Cartographer mapping
    (Section~\ref{sec:cartographer}) added on top of Phase~1,
    confirmed to require no odometry extension at all given the
    \texttt{use\_odometry: false} configuration described in
    Section~\ref{sec:cartographer}. A real, growing occupancy grid was
    built by driving the vehicle around a bounded space under
    teleoperation and saved via \texttt{nav2\_map\_server}'s
    \texttt{map\_saver\_cli} -- this saved map is the one used for all
    subsequent localisation and navigation testing.
  \item[Phase 3 -- Navigation.] The full Nav2 stack (AMCL, planning,
    MPPI control) added on top of Phase~2's saved map, requiring the
    wheel-odometry extension of Section~\ref{sec:onboard-bridge-odom}
    (deferred from Phase~2, where it had not been necessary) and the
    steering-trim (Section~\ref{sec:steering-trim}) and throttle
    calibration (Section~\ref{sec:throttle-calibration}) work, since
    dead-reckoning odometry and AMCL localisation are only as accurate
    as the vehicle model producing them.
\end{description}

The full stack's initial bring-up -- motor bridge, LiDAR node, relay
node, and the ported Nav2 launch configuration -- came up cleanly on
the first attempt against the saved Phase~2 map: both Nav2 lifecycle
managers reported their managed nodes active, all seven configured
MPPI critics (the stock six plus \texttt{EarlyCommitCritic}) loaded
matching the intended configuration, and the costmap auto-sized
correctly to the saved map's dimensions, with no crash loop. A real
navigation goal was deliberately not sent during this first bring-up,
since dead-reckoning odometry -- and therefore everything built on
top of it, including AMCL and planning -- is only as trustworthy as
the still-unresolved steering-trim calibration at that point in the
project's timeline; that calibration (Section~\ref{sec:steering-trim})
was completed before the first real goal was attempted.

\section{A K-Turn Freeze at Reversal Cusps: Diagnosis and a Custom Smoother}
\label{sec:cusp-freeze}

Beyond the display-latency and calibration issues of the preceding
sections, real-hardware testing in tight spaces surfaced a distinct
and more serious navigation problem: a goal requiring the vehicle to
turn around produced, in the operator's words, ``a chaotic curve, then
reverse with opposite curve, then forward curve'' -- a visibly
hesitant, multi-second freeze at the direction change -- rather than a
clean, continuous K-turn. This section documents the diagnosis,
several fix attempts that failed instructively, and a custom smoother
plugin that resolved it, validated in both simulation and on the
physical vehicle.

\subsection{Reproduction}
\label{sec:cusp-repro}

To investigate safely and repeatably, a narrow dead-end corridor
(approximately \SI{2.4}{\meter} wide, \SI{2}{\meter} deep) was added to
the simulation map by re-mapping the Gazebo world of
Chapter~\ref{ch:sim-stack} and saving a fresh occupancy grid via
\texttt{map\_saver\_cli}, then sending a \texttt{NavigateToPose} goal
deep inside it with a reversed final heading -- a target that forces a
K-turn in a space too tight to complete it in one continuous arc.
Two logging tools were built for this investigation:
one recording every \texttt{/plan} message (not only the first, so
that every replan issued during a recovery cycle gets its own
sequence-numbered trace), \texttt{/odom}, and \texttt{/cmd\_vel} to
timestamped CSVs, and a companion analysis script computing total
manoeuvre duration, position-stagnation windows (freeze detection),
and true steering-direction-reversal counts directly from those CSVs.

\subsection{Diagnosis: The Planner Is Correct, the Controller Is Not}
\label{sec:cusp-diagnosis}

Reconstructing the vehicle's true chassis heading through the
manoeuvre (motion direction minus \ang{180} after the cusp, verified
against a independent bicycle-model simulation of a three-point turn)
confirmed that \texttt{SmacPlannerHybrid} was producing a kinematically
correct global plan: heading rotates smoothly and monotonically
through the single Reeds--Shepp cusp the plan contains, exactly the
``flip steering side to continue rotating the same way'' geometry a
real K-turn requires. The freeze is not a planning defect at all --
it is entirely in \emph{execution}.

The root cause lies in \texttt{nav2\_mppi\_controller}'s path-handling
logic (\texttt{path\_handler.cpp}, Humble branch): MPPI is only ever
shown the global plan truncated up to the cusp pose until a function
called \texttt{isWithinInversionTolerances()} reports that the vehicle
has reached that pose within \emph{both} a position tolerance and a
heading tolerance simultaneously -- until then, none of the path
beyond the cusp exists as far as the controller's own path-following
critics are concerned, so nothing in the cost function encourages
progress through it. This is structurally the same failure mode as a
final-goal reorientation freeze found and fixed earlier in this
project's tuning history, outside the scope of this thesis's own
narrative: a goal requiring a large final heading change would stall
in the same way, for the same reason, and was resolved simply by
raising \texttt{GoalAngleCritic}'s cost weight (from 3.0 to 10.0) so
the controller is pushed to complete the final heading precisely. No
equivalent encouragement exists for an \emph{intermediate} cusp pose,
which is exactly where this freeze occurs.

\subsection{Three Plausible Fixes, All Reverted}
\label{sec:cusp-fixes-failed}

Three direct attempts to relax this gating condition were each tested
against the reproduction and reverted, all pointing to the same
underlying lesson:

\begin{itemize}
  \item \textbf{Widening \texttt{isWithinInversionTolerances()}'s own
    position and heading tolerances} (from \SI{0.2}{\meter}/\SI{0.4}{\radian}
    to \SI{0.4}{\meter}/\SI{0.7}{\radian}) made things worse: the
    controller now reported the cusp ``reached'' too loosely, and the
    resulting premature reveal of the post-cusp path triggered
    six or more \texttt{progress\_checker} aborts before the vehicle
    had even physically reached the cusp.
  \item \textbf{Faster stuck-detection} (lowering
    \texttt{progress\_checker}'s movement-time allowance from
    \SI{30}{\second} to \SI{8}{\second}) fired aborts every
    \SIrange{8}{13}{\second} -- faster than the manoeuvre could
    genuinely complete -- and once triggered a \texttt{BackUp}
    recovery action that itself failed outright for lack of room in
    the corridor.
  \item \textbf{A first-version path-smoother} that re-projected the
    path onto a straight line for a fixed distance after each cusp,
    then jumped instantly back to the original curve, performed worse
    still (never completing within a \SI{108}{\second} window): the
    plugin's own logic worked exactly as designed, but the hard
    discontinuity it introduced at the blend boundary was itself a new
    obstacle for the controller.
\end{itemize}

All three attempts shared a common thread that became the guiding
principle for the eventual fix: \emph{the freeze needs real,
uninterrupted time and space to resolve on its own; interrupting it,
however that interruption is triggered, consistently made the
manoeuvre worse, never better.}

\subsection{A Reverse-Only, Gradually Blended Smoother}
\label{sec:cusp-smoother-v2}

A second design, \texttt{CuspStraightenerSmoother}, was built as a
\texttt{nav2\_core::Smoother} plugin -- path post-processing rather
than a controller-side change -- with two deliberate constraints: it
acts \emph{only} on cusps the path enters in reverse (a front-steered
vehicle's steered axle is self-correcting driving forward but
trailing, and jackknife-prone, in reverse, making a reversing cusp the
harder sub-case; forward-entering cusps and all ordinary forward
segments are left completely untouched), and it blends
\emph{gradually} rather than cutting over instantly, tapering linearly
from the straightened projection at the cusp itself back to the
original planned curve over a configurable \texttt{straight\_distance}.

\paragraph{A mutation-order bug, and the fix.} Directly testing this
design against harder, multi-cusp replans (plausible whenever a path is
replanned from a robot pose already mid-manoeuvre, e.g.\ during a
recovery cycle) surfaced a genuine implementation bug: the original
\texttt{smooth()} detected every cusp once up front, then processed
each one \emph{in place} on the very path array it was simultaneously
rewriting. When two cusps fell close together, an earlier cusp's blend
could overwrite position and heading data that a nearby later cusp
still needed to read in order to classify itself and compute its own
straight-line target -- a self-inflicted data race within a single,
single-threaded function. The fix, shown in
Listing~\ref{lst:cusp-detect}, snapshots the incoming path into a
read-only \texttt{original} vector before any modification begins;
every geometric decision the function makes -- cusp detection,
forward/reverse classification, and every blend target -- reads only
from this untouched snapshot, and writes go exclusively to a separate
output path being built alongside it.

\begin{lstlisting}[language=C++,caption={Cusp detection and reverse/forward classification, reading only from a read-only snapshot (\texttt{cusp\_straightener\_smoother.cpp}).},label={lst:cusp-detect}]
const std::vector<geometry_msgs::msg::PoseStamped> original = path.poses;

std::vector<size_t> cusp_indices;
for (size_t i = 1; i + 1 < n; ++i) {
  const double d1x = original[i].pose.position.x - original[i - 1].pose.position.x;
  const double d1y = original[i].pose.position.y - original[i - 1].pose.position.y;
  const double d2x = original[i + 1].pose.position.x - original[i].pose.position.x;
  const double d2y = original[i + 1].pose.position.y - original[i].pose.position.y;
  const double len1 = std::hypot(d1x, d1y);
  const double len2 = std::hypot(d2x, d2y);
  if (len1 < 1e-6 || len2 < 1e-6) { continue; }
  const double dot = (d1x * d2x + d1y * d2y) / (len1 * len2);
  if (dot < 0.0) { cusp_indices.push_back(i); }  // direction reverses by > 90 deg
}
// Classify each cusp (forward- vs reverse-entering) from `original` only,
// independent of anything the output-building loop below has written so far.
for (const size_t cusp_i : cusp_indices) {
  const double cusp_yaw = tf2::getYaw(original[cusp_i].pose.orientation);
  const double dx = original[cusp_i + 1].pose.position.x - original[cusp_i].pose.position.x;
  const double dy = original[cusp_i + 1].pose.position.y - original[cusp_i].pose.position.y;
  const double along = dx * std::cos(cusp_yaw) + dy * std::sin(cusp_yaw);
  cusps.push_back({cusp_i, along < 0.0, along});  // along < 0 => reverse-entering
}
\end{lstlisting}

Directly re-tested afterward against a genuine multi-cusp replan (a
47-point path with three cusps, two of them at adjacent indices), the
fixed version classified and handled both correctly -- exactly the
scenario the original bug would have corrupted.

Table~\ref{tab:cusp-trials} compares three trials each of the
untouched baseline, the buggy v2 smoother, and the fixed v2 smoother,
on the identical corridor reproduction goal with a fresh stack
relaunch before every trial.

\begin{table}[htbp]
  \centering
  \caption{K-turn manoeuvre duration and recovery count, corridor reproduction goal (simulation).}
  \label{tab:cusp-trials}
  \begin{tabular}{@{}lccc@{}}
    \toprule
    Configuration & Trial durations (recoveries) & Mean & Worst case \\
    \midrule
    Baseline (no smoother) & 81.3s(1), 81.1s(1), 168.1s(4) & 110.2s / 2.0 & 168.1s \\
    v2 smoother (buggy) & 96.6s(2), 101.4s(2), 432.8s(12) & 210.3s / 5.3 & 432.8s \\
    v2 smoother (fixed) & 59.3s(0), 97.1s(2), 103.7s(2) & 86.7s / 1.3 & 103.7s \\
    \bottomrule
  \end{tabular}
\end{table}

Run-to-run variance is genuinely high in every configuration --
consistent with the freeze being a stochastic MPPI sampling effect
rather than a deterministic planning defect -- but the fixed smoother
clearly outperforms the untouched baseline on every measure, while the
buggy version is markedly worse than doing nothing at all, underlining
how much the mutation-order bug alone was costing.

\subsection{``Extending the Corner'': A Further, Larger Improvement}
\label{sec:cusp-extend}

A further refinement addressed the residual freeze directly at its
source: rather than requiring the path's incoming and outgoing legs to
meet at one precise geometric vertex -- which
\texttt{isWithinInversionTolerances()} must still converge on exactly,
even with the gradual blend of Section~\ref{sec:cusp-smoother-v2} in
place -- a new \texttt{extend\_distance} parameter pushes the cusp
point itself further along the incoming leg's own heading before the
reverse blend begins, so the incoming and outgoing legs deliberately
overlap rather than meeting at a single tight point, giving the
controller geometric margin around the transition instead of an exact
target. Listing~\ref{lst:cusp-extend} shows the construction: the
extension point is inserted into the output path, and the existing
distance-based blend is re-anchored to start from it rather than from
the original cusp position.

\begin{lstlisting}[language=C++,caption={Pushing the cusp outward before blending, so incoming/outgoing legs overlap (\texttt{cusp\_straightener\_smoother.cpp}).},label={lst:cusp-extend}]
const double cusp_yaw = tf2::getYaw(original[i].pose.orientation);
const double ex = original[i].pose.position.x + std::cos(cusp_yaw) * extend_distance_;
const double ey = original[i].pose.position.y + std::sin(cusp_yaw) * extend_distance_;

geometry_msgs::msg::PoseStamped ext_pose = original[i];
ext_pose.pose.position.x = ex;
ext_pose.pose.position.y = ey;
output.push_back(ext_pose);
// The existing straight_distance_ blend (unchanged logic) now runs from
// (ex, ey) as its origin, not from the original cusp position - the reverse
// leg must travel back past this pushed-out point to rejoin the downstream
// curve, giving the controller a margin instead of one exact vertex.
\end{lstlisting}

With \texttt{extend\_distance} set to \SI{0.15}{\meter} alongside the
existing \SI{0.21}{\meter} \texttt{straight\_distance}, three further
trials produced \SI{45.2}{\second}, \SI{38.3}{\second}, and
\SI{48.1}{\second} -- a mean of \SI{43.9}{\second} with \textbf{zero
recoveries in every trial}, roughly halving the already-improved
fixed-v2 mean of \SI{86.7}{\second} on top of its improvement over the
\SI{110.2}{\second} baseline. This is the final, locked-in
configuration.

\subsection{Real-Hardware Validation}
\label{sec:cusp-hw-validation}

The fixed, corner-extending smoother was tested on the physical QCar
with two goals sent manually by the operator, each requiring a genuine
reversal manoeuvre: a \SI{5.15}{\meter} traverse completed in
\SI{50.1}{\second} with zero recoveries across seven replans (cusp
counts 1, 1, 1, 1, 8, 8, 0), and a second, longer traverse completed in
\SI{96.8}{\second}, also with zero recoveries, across eight replans
(cusp counts 5, 5, 5, 5, 9, 10, 5, 3). Both goals succeeded outright --
the strongest available confirmation that the fix generalises beyond
the simulated corridor reproduction it was developed against.

Both goals also produced noticeably higher cusp counts on individual
replans (up to 8--10) than any simulation trial had shown (maximum
approximately 3) -- an open, unexplained gap between simulated and
real replanning behaviour, plausibly \texttt{SmacPlannerHybrid}
responding to more awkward real goal/environment geometry than the
simulated corridor exercised, but not confirmed. It did not prevent
either goal from succeeding, so it is recorded here as a genuine open
item for future work (Section~\ref{sec:prospects}) rather than a
present limitation of the fix.

One apparent timing anomaly in this trace was investigated and
correctly ruled out rather than mistaken for a bug: the first several
\texttt{smooth()} calls on the first goal landed within roughly one
second of each other, which initially looked like an unthrottled
replan burst. Re-reading the behaviour tree configuration showed this
is by design: \texttt{IsPathValid} is already rate-limited to
\SI{1}{\hertz} by an existing \texttt{RateController}, and a
\texttt{RateController} fires immediately on its first tick before
settling into its steady interval -- exactly the pattern observed,
once the also-unconditional, once-only \texttt{InitialComputePathToPose}
call is accounted for separately. Catching this before implementing a
second, redundant rate limiter avoided a change that would have had no
effect at all.

\subsection{Continued Iteration: A Re-Extension Bug and a Curvature-Continuing Refinement}
\label{sec:cusp-continued-iteration}

Hardware testing continued beyond the validation above, and surfaced a
further, genuine bug in the corner-extension mechanism itself, plus a
refinement to it -- both developed and sim-tested after the
hardware-validated result of Section~\ref{sec:cusp-hw-validation}, and
not yet re-confirmed on the physical vehicle at the time of writing;
they are reported here with that status made explicit, consistent with
this thesis's practice of not presenting a provisional result as final.

\paragraph{Bug: the extension re-applies to the same cusp on every replan.} Because Nav2 replans roughly once per second while a goal is
active, and each replan's path starts fresh from the robot's current
pose, the \emph{same} physical K-turn cusp reappears near the front of
the replanned path on every single replan as the robot approaches and
then executes it -- and, since the smoother has no notion of ``already
handled this one,'' it re-applied the full extend-and-straighten
treatment to that same cusp every time. Direct log analysis of a real
hardware run confirmed this concretely: one K-turn's cusp received the
treatment on eight separate replans as the robot closed in on it, its
index in the (shrinking) replanned path falling from 94 down to 1 --
by the final replan, the robot was already essentially at the cusp,
mid-manoeuvre, and the smoother was still reshaping the path directly
underneath it rather than leaving an already-committed manoeuvre
alone. This is, functionally, a milder recurrence of the same
``interrupting an in-progress manoeuvre makes it worse'' pattern
identified repeatedly elsewhere in this investigation
(Section~\ref{sec:cusp-fixes-failed}), this time self-inflicted by the
smoother rather than by an external replan trigger.

The fix adds a \texttt{min\_lead\_distance} parameter (\SI{0.3}{\meter}):
the arc length from the replanned path's start to each candidate cusp
is computed, and the extend-and-straighten treatment is only applied
if the cusp is still at least this far ahead of the robot's current
position, leaving an already-reached or already-underway cusp
untouched rather than reshaping it again. Two sim trials on the
standard K-turn reproduction goal after this fix showed no regression,
though the specific skip condition was not directly observed firing in
either trial (both completed before any replan's cusp closed to within
the threshold) -- correct by code review and by the logged parameter
value at start-up, but not yet directly confirmed exercising its skip
branch. This is recorded as an open verification item, not a claimed
result.

\paragraph{Refinement: continuing the incoming curve rather than cutting straight.} Separately, the extension itself was upgraded from a
straight-line push-out (Listing~\ref{lst:cusp-extend}, hardware-validated
above) to one that continues the incoming leg's own curvature. The
motivation is direct: cutting straight right at the corner, even when
the vehicle is still visibly curving on its approach to the cusp, asks
the vehicle to straighten abruptly at exactly the point it should be
easing into the manoeuvre most naturally. The refined version
estimates the incoming leg's curvature at the cusp from its own recent
yaw change, $\kappa = \Delta\theta/\Delta s$ over the last incoming
segment, and generates the extension as a short constant-curvature arc
with that same curvature rather than a straight line -- built from
several closely spaced sub-poses (\SI{0.03}{\meter} steps) so the
curve is genuinely present in the path's own geometry, not merely
implied by a single endpoint. It degrades automatically to the
original straight-line behaviour when the incoming leg is already
straight ($\kappa \approx 0$), so it is a strict generalisation of the
hardware-validated mechanism rather than a replacement with different
behaviour on straight approaches.

\begin{lstlisting}[language=C++,caption={Continuing the incoming leg's own curvature into the extension, rather than cutting straight (\texttt{cusp\_straightener\_smoother.cpp}).},label={lst:cusp-curvature}]
// kappa estimated from the incoming leg's own yaw change (dyaw/ds) at the cusp - 0 if the
// incoming leg was already straight, giving the original straight-line behaviour unchanged.
constexpr double kExtensionStep = 0.03;
const int extension_steps = std::max(1, static_cast<int>(std::round(extend_distance_ / kExtensionStep)));
const double step = extend_distance_ / extension_steps;
double ex = cx, ey = cy, ext_yaw = cusp_yaw;
for (int s = 1; s <= extension_steps; ++s) {
  const double arc = step * s;
  double px, py;
  if (std::abs(kappa) > 1e-6) {
    const double radius = 1.0 / kappa;
    px = cx + radius * (std::sin(cusp_yaw + kappa * arc) - std::sin(cusp_yaw));
    py = cy - radius * (std::cos(cusp_yaw + kappa * arc) - std::cos(cusp_yaw));
  } else {
    px = cx + std::cos(cusp_yaw) * arc;   // kappa ~ 0: degrades to the straight-line case
    py = cy + std::sin(cusp_yaw) * arc;
  }
  ext_yaw = cusp_yaw + kappa * arc;
  // ... each (px, py, ext_yaw) is appended as its own pose, making the arc explicit.
}
\end{lstlisting}

A single sim smoke test on the standard reproduction goal completed
successfully ($\approx$\SI{99}{\second}, within the range of prior
trials) with the new curvature-extension code path exercised across
eight reverse-cusp events over the run's several replans, with no
numerical failure. This is reported as a working, sim-verified
refinement, not yet subjected to the same multi-trial statistical
comparison or hardware re-validation as the straight-line version it
replaces -- a concrete, already-identified next step
(Section~\ref{sec:prospects}).

\subsection{A Related Planner-Side Finding: Near-Wall Goals and a Tolerance Trade-off}
\label{sec:near-wall-tolerance}

A separate, planner-level finding surfaced during the same period of
continued tuning, also directly relevant to operating in tight, wall-adjacent
spaces: a goal placed only $\approx$\SI{0.20}{\meter} from the nearest
mapped wall -- inside the footprint clearance the vehicle's own
\SI{0.44}{\meter}$\times$\SI{0.18}{\meter} collision footprint requires
-- failed outright on every one of six recovery retries, each attempt
logging \texttt{SmacPlannerHybrid} exhausting its entire search
iteration budget without ever finding a collision-free plan. This is a
goal-placement issue rather than a code defect: no footprint- or
inflation-related configuration was implicated, and the practical
guidance that follows is simply to keep any commanded goal at least
\SIrange{0.3}{0.4}{\meter} clear of a mapped obstacle, given this
vehicle's footprint.

A tempting-looking fix was tried and deliberately reverted, and the
reason why is itself a useful finding. \texttt{SmacPlannerHybrid}'s
own search loop, confirmed directly against its source, backtracks to
the closest node it reached within a configurable \texttt{tolerance}
radius rather than failing outright whenever the exact goal proves
infeasible -- so widening \texttt{tolerance} (from \SI{0.5}{\meter} to
\SI{1.0}{\meter}) did immediately rescue the failing near-wall goal in
a live test. However, it introduced a worse, silent failure mode in
its place: the substituted pose the planner backtracks to, not the
originally requested goal, becomes the plan's actual endpoint, and
nothing downstream in the standard \texttt{NavigateToPose} behaviour
tree ever rechecks the final pose against the original request --
so \emph{any} marginally infeasible goal, not only the intended
near-wall case, could now report \texttt{Goal succeeded} while the
vehicle sat up to a full \SI{1.0}{\meter} away from, and potentially
facing the wrong direction relative to, where it was actually asked to
go, with no indication in the result that a substitution had occurred
at all. An honest, visible failure is preferable to a silently
mis-reported success, particularly for a system whose results this
thesis reports as measured fact -- \texttt{tolerance} was reverted to
its original \SI{0.5}{\meter}, accepting that the specific near-wall
goal will fail outright again, since the correct fix for a genuine
target at that location is to ensure real clearance there, not to
relax what counts as ``success'' for every goal system-wide. This
mirrors, at the planner level, the same safety-over-convenience
judgement Section~\ref{sec:active-braking} applies at the control
level: a fix that resolves the reported symptom by quietly weakening a
correctness guarantee elsewhere is not an acceptable trade, even under
time pressure.

\subsection{Two Follow-On Attempts, Reverted}
\label{sec:cusp-followon-reverted}

Two further improvements were implemented and hardware-tested the same
day, and both were reverted -- included here because a reverted
attempt with a documented, evidence-based reason is exactly the kind
of result this thesis's methodology (Section~\ref{sec:methodology-overview})
treats as worth reporting, not discarding.

\texttt{IsPathValidDebouncedCondition}, a custom behaviour-tree
condition requiring two consecutive failed validity checks (rather than
stock \texttt{IsPathValid}'s single check) before triggering a replan,
improved simulation trials further still (\SIlist{38.2;38.3;40.4}{\second},
mean \SI{39.0}{\second}, zero recoveries) by absorbing occasional
single-tick sensor or transform noise that would otherwise trigger an
unnecessary replan. On real hardware, however, a subsequent run on a
freshly scanned map showed the vehicle never executing its reverse
segment at all: the local costmap reported the path invalid on
\emph{every} check, not intermittently, so the debounced condition
correctly (and unhelpfully) confirmed a persistent rather than
transient invalidity every \SIrange{1.5}{2.6}{\second}, aborting the
in-progress manoeuvre each time -- the same ``interrupting a manoeuvre
makes it worse'' lesson from Section~\ref{sec:cusp-fixes-failed},
under a new trigger. The most likely root cause, not fully confirmed
before the revert, is an unmeasured \texttt{qcar\_clock\_offset}
(Section~\ref{sec:clock-offset}) on the relay for that session -- a
fresh measurement immediately afterward found a real
$\approx$\SI{0.15}{\second} offset. A companion \texttt{max\_cusps}
rejection cap was reverted alongside it to return to one single
known-good state rather than debug two independent changes under time
pressure, though it was not itself implicated (every replan in the
failing run came back at exactly two cusps).

A custom \texttt{PathDeviationCritic} -- pulling sampled trajectories
toward a point further along the path once the vehicle drifts more
than \SI{0.3}{\meter} from it, rather than toward the nearest path
point -- was tested for three trials
(\SIlist{112.8;71.7;38.4}{\second}, mean \SI{74.3}{\second}, zero
recoveries) against the corner-extending smoother's own
\SI{43.9}{\second} mean, and reverted: both the mean and the
run-to-run variance were worse, plausibly because the new critic's
pull competed directly against the stock \texttt{PathAlignCritic}
rather than complementing it.

\subsection{A Related Operational Caution: Goal Preemption}
\label{sec:goal-preemption-caution}

A separate mechanics finding from this investigation, worth recording
as an operational rule rather than a fix: if a goal is deep into a
failing/retrying cycle (for example, the planner repeatedly failing to
find a plan, with the behaviour tree's recovery retries nearly
exhausted) and a new goal is sent while the old one is still active, the
new goal does not reliably receive a fresh attempt -- it was observed
to fail almost instantly, with no real planning attempt of its own,
rather than superseding the old goal cleanly. The practical rule is to
let a stuck goal finish (fail out) or explicitly cancel it before
sending a correction, rather than assuming a new goal always starts
with a clean slate.

\section{Diagnosing an Apparent Display-Localisation Artefact}
\label{sec:display-lag}

After the throttle calibration of Section~\ref{sec:final-calibration}
(but before the buffering-defect fix of Section~\ref{sec:readmode-bug}),
a live Nav2 goal test surfaced a new symptom, described by the
vehicle's operator in detail: when driving toward a goal inside an
enclosed space, the entire mapped rectangle appeared to shift in
RViz, with the LiDAR scan keeping its internal shape but dragging
behind the robot's true motion before ``snapping'' back into
alignment -- and the same shift-then-correct pattern recurred, after
a delay, once the goal was reported reached.

\subsection{Ruling Out Simpler Causes}
\label{sec:display-lag-rule-out}

Several candidate causes were checked directly and eliminated:
\texttt{qcar\_clock\_offset} was re-measured fresh and re-applied
(with one procedural caution found along the way -- a relay launch
command given during this session had omitted the
\texttt{qcar\_clock\_offset} argument entirely, silently defaulting
it to zero and briefly reintroducing exactly the bug fixed in
Section~\ref{sec:clock-offset}, worth checking for explicitly before
suspecting the offset value itself); \texttt{map}
$\rightarrow$ \texttt{odom} TF was captured directly during active
driving and updated in small, smooth, incremental steps rather than
sudden jumps; the RViz \texttt{Fixed Frame} configuration was
confirmed correct (\texttt{map}); \texttt{base} $\rightarrow$
\texttt{lidar} TF is static and therefore not a timing bottleneck;
local CPU load was checked and found unremarkable (RViz
$\approx$10\%, the Nav2 container $\approx$29\%); network round-trip
latency to the QCar was \SIrange{1}{6}{\milli\second}, far too small
to explain a multi-second visible effect; and direct
receipt-versus-capture latency on \texttt{/odom} and \texttt{/scan}
measured $\approx$\SIlist{0.18;0.20}{\second} respectively, small
and tightly clustered, ruling out network or relay transit as the
source of a large visible lag.

\subsection{Instrumentation: \texttt{monitor\_goal\_timing.py}}
\label{sec:monitor-goal-timing}

A dedicated instrumentation tool was built to timestamp a single Nav2
goal end to end: goal-sent, real motion start/stop (from
\texttt{/odom}), RViz-displayed motion start/stop (from the
\texttt{map} $\rightarrow$ \texttt{base} TF), and goal-reached. Two
implementation defects were found and fixed while building this tool,
both worth recording as general lessons for any future action-status
monitoring code on ROS~2:

\begin{enumerate}
  \item \textbf{A QoS durability mismatch.} The actual publisher QoS
    for \texttt{/navigate\_to\_pose/\_action/status} is
    \texttt{RELIABLE} + \texttt{TRANSIENT\_LOCAL}. \texttt{rclpy}'s
    default subscription QoS, constructed by passing a plain integer
    depth, leaves durability at \texttt{VOLATILE} -- technically
    QoS-compatible by specification (a subscriber may request less
    than a publisher offers), but this delivered \emph{zero} messages
    in practice, confirmed twice on real hardware. The fix required
    constructing an explicit \texttt{QoSProfile} matching the
    publisher exactly, a reminder that formal QoS compatibility rules
    do not guarantee the practical behaviour a developer expects
    (Section~\ref{sec:ros2-model}). This is exactly the
    \texttt{TRANSIENT\_LOCAL}-publisher / \texttt{VOLATILE}-subscriber
    case Table~\ref{tab:qos-compatibility} marks as formally
    compatible: durability governs whether \emph{already-published}
    samples are replayed to a newly matched subscriber, and a
    low-frequency topic like action status may not publish a fresh
    sample again for a long time, so a \texttt{VOLATILE} subscriber
    that missed the replay is left waiting with no error raised on
    either side.
  \item \textbf{Goal-detection logic.} The tool originally waited for
    status code \texttt{ACCEPTED} specifically to mark ``goal sent'',
    but a goal can transition directly from absent to
    \texttt{EXECUTING} between two consecutive status publishes,
    skipping an independently observable \texttt{ACCEPTED} update
    entirely -- this topic's \texttt{TRANSIENT\_LOCAL} history appears
    to retain only the latest status per goal, not the full transition
    sequence. The fix tracks goal UUIDs instead: the first appearance
    of a previously unseen \texttt{goal\_id} (after seeding a known-ID
    set from whatever the topic already retained at monitor startup)
    marks ``goal sent'', regardless of which status value it first
    carries.
\end{enumerate}

Listing~\ref{lst:status-qos} shows the explicit \texttt{QoSProfile}
that fixed the first defect -- matching \texttt{bt\_navigator}'s
publisher exactly rather than relying on the default subscription QoS
constructed from a plain integer depth.

\begin{lstlisting}[language=Python,caption={Explicit QoS profile matching \texttt{bt\_navigator}'s action-status publisher (\texttt{scripts/monitor\_goal\_timing.py}).},label={lst:status-qos}]
STATUS_QOS = QoSProfile(
    reliability=QoSReliabilityPolicy.RELIABLE,
    durability=QoSDurabilityPolicy.TRANSIENT_LOCAL,
    history=QoSHistoryPolicy.KEEP_LAST,
    depth=10,
)
...
self.create_subscription(GoalStatusArray, '/navigate_to_pose/_action/status',
                          self.status_cb, STATUS_QOS)
\end{lstlisting}

Listing~\ref{lst:goal-uuid-tracking} shows the second fix: rather than
waiting for a specific status code, the callback seeds a
\texttt{known\_goal\_ids} set from whatever the topic already retains
at startup (so pre-existing, already-completed goals are not
mistaken for a new one), then marks ``goal sent'' the first time a
previously unseen goal UUID appears in any status, independent of
which status value that UUID first carries.

\begin{lstlisting}[language=Python,caption={Goal-UUID--based ``goal sent'' detection (\texttt{scripts/monitor\_goal\_timing.py}).},label={lst:goal-uuid-tracking}]
if not self.known_goal_ids:
    for s in msg.status_list:
        self.known_goal_ids.add(bytes(s.goal_info.goal_id.uuid))
    self.get_logger().info(
        'seeded %d pre-existing goal id(s), waiting for a new one'
        % len(self.known_goal_ids))
    return

for s in msg.status_list:
    gid = bytes(s.goal_info.goal_id.uuid)
    if gid not in self.known_goal_ids:
        self.known_goal_ids.add(gid)
        if self.goal_sent_t is None:
            self.goal_sent_t = now
            self.tracked_goal_id = gid
\end{lstlisting}

An explicit discovery-wait (polling \texttt{get\_publishers\_info\_by\_topic}
plus a settle margin) was also required before treating the monitor
as ready to observe a goal -- a subscription created and immediately
assumed matched can miss a fast goal's status entirely if DDS
discovery with the separate \texttt{bt\_navigator} process has not
yet completed.

\subsection{Measured Result}
\label{sec:display-lag-result}

Table~\ref{tab:display-lag} presents one complete goal traced with
this tool.

\begin{table}[htbp]
  \centering
  \caption{Timing trace of one Nav2 goal (\texttt{monitor\_goal\_timing.py}).}
  \label{tab:display-lag}
  \begin{tabular}{@{}lll@{}}
    \toprule
    Event & Offset & Pose \\
    \midrule
    Goal sent & -- & -- \\
    Real start (odom) & +2.145\,s after goal sent & (3.344, 1.674) odom frame \\
    RViz start (map$\to$base) & +2.705\,s after goal sent & (0.033, $-$0.031) map frame \\
    Goal reached & +6.660\,s after goal sent & -- \\
    RViz stop & +2.094\,s after goal reached & (1.257, $-$0.024) \\
    Real stop (odom) & +2.238\,s after goal reached & -- \\
    \bottomrule
  \end{tabular}
\end{table}

This gives a start gap (RViz minus real) of $+\SI{0.559}{\second}$
and a stop gap of $-\SI{0.144}{\second}$ (RViz reporting ``stopped''
slightly \emph{before} the real vehicle fully settled -- small, and
plausibly within the detection threshold's own noise rather than a
separate effect). Note that the real-start offset of
\SI{2.145}{\second} already includes the static-friction stiction
delay characterised in Section~\ref{sec:deadband-discovery}, so it is
not, by itself, a new timing number.

\subsection{Interpretation, and a Necessary Caveat}
\label{sec:display-lag-interpretation}

This measurement supports two separate, unrelated conclusions. The
$\approx$\SI{0.56}{\second} start-of-goal gap is a modest, constant
pipeline latency -- network transit, AMCL's scan-matching cycle, and
RViz's own subscribe/render overhead stacking up -- between a real
event and its appearance on screen; it is not a TF-jump defect (ruled
out in Section~\ref{sec:display-lag-rule-out}), and it is not a real
navigation-performance problem, since Nav2's own controller consumes
\texttt{/odom}/\texttt{/amcl\_pose}/\texttt{/tf} directly at the
source, without RViz's additional rendering overhead -- the robot's
actual control decisions are working with data less stale than what a
human operator sees on screen. The end-of-goal gap, by contrast, is
the same stop-latency/coast phenomenon already characterised in
Chapter~\ref{ch:calibration}: the real vehicle continuing to move for
\SI{2.238}{\second} after ``goal reached'' matches the pre-fix
stop-latency figures in Table~\ref{tab:calibration-validation}'s
predecessor almost exactly, confirmed here via a completely
independent measurement method (a live Nav2 goal rather than the
open-loop \texttt{drive\_and\_log.py} test).

The necessary caveat, and the reason this section is presented in
detail rather than as a closed result: \textbf{this entire
measurement was taken while the \texttt{readMode=1} buffering defect
of Section~\ref{sec:readmode-bug} was still active} in
\texttt{qcar\_bridge.py} -- it was fixed the following day. Given that
defect's confirmed effect on every encoder-derived quantity
(Table~\ref{tab:readmode-fix}), the \SI{2.238}{\second} end-of-goal
figure here is almost certainly, in large part, the same artefact
rather than a distinct navigation-specific finding -- its magnitude
matches the buffering defect's characterised effect too closely to be
coincidental. The \SI{0.559}{\second} start-gap figure is a different
measurement (\texttt{map}$\to$\texttt{base} TF timing rather than
\texttt{/odom} velocity thresholding) and is less obviously the same
artefact, but since it is also downstream of the same buffered
\texttt{/odom} stream, it cannot be treated as a clean, artefact-free
measurement of pure network/AMCL/RViz pipeline latency either. The
diagnostic methodology and tooling built in this section -- the
QoS/goal-detection fixes, and the systematic elimination of TF-jump,
Fixed-Frame, and network-latency causes -- remain valid and reusable;
only the specific numbers in Table~\ref{tab:display-lag} should be
treated as provisional pending a fresh \texttt{monitor\_goal\_timing.py}
run with the \texttt{readMode=0} fix in place. % TODO: re-run and report

\section{Final Validation}
\label{sec:final-validation}

Following the \texttt{readMode=0} fix (Section~\ref{sec:readmode-bug}),
a live Nav2 goal was re-run on the physical vehicle -- full onboard
bridge, LiDAR node, and relay node, with AMCL localised against the
existing saved map. The vehicle's operator's direct, contemporaneous
assessment was unambiguous: navigation was ``working almost
perfectly'' -- a strong qualitative result, and one consistent with
the buffering defect having been the dominant cause of the
``skeleton of the map is moving'' symptom that opened
Section~\ref{sec:display-lag}, rather than merely a contributing
factor to it. \texttt{monitor\_goal\_timing.py} was not re-run for
this specific goal, so the precise post-fix start-gap and
end-of-goal figures corresponding to Table~\ref{tab:display-lag}
are not available as quantitative numbers in this thesis. % TODO:
% capture these numbers with a fresh lab session if at all possible -
% this is the single most valuable remaining data point for Ch. 7.

\section{Summary}
\label{sec:nav-integration-summary}

This chapter closes the sim-to-real migration begun in
Chapter~\ref{ch:sim-stack}: the physical QCar, connected through the
bridge of Chapter~\ref{ch:hardware-bridge} and calibrated per
Chapter~\ref{ch:calibration}, was shown to complete real, closed-loop
Nav2 navigation goals, including an Ackermann-appropriate custom MPPI
critic and phased integration that isolated each subsystem's own
correctness before combining it with the next. Section~\ref{sec:cusp-freeze}'s
custom smoother resolved a genuine execution-side freeze at reversal
cusps -- reducing a representative K-turn manoeuvre's simulated mean
duration from \SI{110.2}{\second} with recoveries to \SI{43.9}{\second}
with none, and validated directly on the physical vehicle across two
independent goals with zero recoveries in either -- while its two
reverted follow-on attempts, and the still-open gap between simulated
and real replan complexity, are reported alongside it rather than
omitted. The one remaining significant open item is quantitative
rather than functional: this chapter's display-lag investigation
(Section~\ref{sec:display-lag}) was measured under a bridge
configuration later found to be affected by the buffering defect of
Section~\ref{sec:readmode-bug}, and a qualitative post-fix confirmation
exists, but a fresh quantitative re-measurement does not.
Chapter~\ref{ch:results} draws together the quantitative results
available across this and the preceding two chapters.
